\begin{acknowledgments}
First, I thank my advisor, Dr.~Ziwen Pan, who gave me the freedom to work across simulation, algorithms, machine learning, error correction and cryptography. Dr.~Pan also guided me to make those threads into one argument. I thank my committee, Dr.~Zihao Zhan, Dr.~Susan Mengel and Dr.~Lu Wei, and the Dean's representative, Dr.~Anne E.~V.~Gorden, for their time, questions and careful reading of this document.

We gratefully thank Dr.~Jan Balewski of the National Energy Research Scientific Computing Center at Lawrence Berkeley National Laboratory for his expertise in running quantum workloads on Perlmutter, which made the large-scale experiments possible, and for the QCrank encoding that underlies several chapters. I thank Alex Khan, Wenshuo Hu, Steven Rayan, Kewen Xiao, Jie Li, Yong Chen, Shabnam Jabeen and Randy Kuang for their collaboration on the papers that form the chapters of this dissertation.

This research used resources of the National Energy Research Scientific Computing Center, a DOE Office of Science User Facility supported by the Office of Science of the U.S. Department of Energy under Contract No.~DE-AC02-05CH11231, using NERSC awards DDR-ERCAP0034486 and DDR-ERCAP0034385. The Texas Tech University High Performance Computing Center also supported this research. The hardware experiments used IBM Quantum services and IonQ systems. The views expressed are my own and do not reflect the official policy or position of IBM, IonQ or their teams. I thank Xanadu for the Amazon Braket credits that funded early trapped-ion runs, and Light Rider Inc.\ for supporting the security work of the final chapters.

Finally, I thank my family and friends for their patience and encouragement throughout this degree.
\end{acknowledgments}
